\documentclass[a4paper, notitlepage, 10pt]{article}
\usepackage{geometry}
\usepackage{titling}
\usepackage{authblk}

% Page Layout Configuration
\geometry{
 a4paper,
 total={170mm,257mm},
 left=20mm,
 top=20mm,
}

% Font and Typography
\usepackage{setspace,relsize}
\usepackage{moreverb}

% Tables, Figures, and Listings
\usepackage{booktabs}
\usepackage{graphicx}
\usepackage{subcaption}
\usepackage{tikz}
\usetikzlibrary{shapes}
\usepackage{wrapfig}
\usepackage{listings}

% Mathematical Packages
\usepackage{amsthm}
\usepackage{amsmath}
\usepackage{amssymb}
\usepackage{mathtools}
\usepackage{bbm}

% Text Formatting
\usepackage{indentfirst}
\usepackage{multirow}
\usepackage{enumitem}
\usepackage{cancel}
\usepackage[normalem]{ulem}

% Hyperlinks and Color
\definecolor{keywords}{RGB}{255,0,90}
\definecolor{comments}{RGB}{0,0,113}
\definecolor{red}{RGB}{160,0,0}
\definecolor{green}{RGB}{0,150,0}

\usepackage{hyperref}
\usepackage{xcolor}
\hypersetup{
    colorlinks=true,
    linkcolor=blue,      % Internal links
    citecolor=blue,      % Citations
    urlcolor=blue        % External hyperlinks
}
\usepackage[square,numbers]{natbib}

% Code Listings Style
\lstset{
    language=Python, 
    basicstyle=\ttfamily\footnotesize, 
    keywordstyle=\color{blue},
    commentstyle=\color{green},
    stringstyle=\color{red},
    showstringspaces=false
}

% Title and Authors
\title{\Large \bfseries Relevant Theorems in Analysis and Probability Theory}
\date{\today}
\author{Lucas Machado Moschen}
\affil{\small Department of Mathematics, Imperial College London}

%%%%%%%%%%%%%%%%%%%%%%%%%%%%%%%%%%%%%%%%%%%%%%%%
%%%%%%%% Mathematical Definitions %%%%%%%%%%%%%%
%%%%%%%%%%%%%%%%%%%%%%%%%%%%%%%%%%%%%%%%%%%%%%%%

\newcommand{\R}{\mathbb{R}}
\newcommand{\bb}[1]{\boldsymbol{#1}}
\newcommand{\Space}{\mathcal{X}}
\newcommand{\C}{\mathbb{C}}
\newcommand{\D}{D}
\newcommand{\B}{\mathcal{B}}

\newcommand{\ev}{\mathbb{E}}
\newcommand{\var}{\operatorname{Var}}
\newcommand{\cor}{\operatorname{Cor}}
\newcommand{\cov}{\operatorname{Cov}}
\newcommand{\almost}{\text{a.e. }}

\newcommand{\ind}{\mathbbm{1}}

% Theorem Environments

\newtheorem{theorem}{Theorem}[section]
\newtheorem{claim}{Claim}[section]
\newtheorem{proposition}{Proposition}[section]
\newtheorem{corollary}{Corollary}[section]
\newtheorem{lemma}{Lemma}[section]

\theoremstyle{definition}
\newtheorem{example}{Example}[section]

\newtheoremstyle{customstyle} % Define a custom style for assumptions
  {10pt}      % Space above
  {10pt}      % Space below
  {\normalfont\setlength{\parindent}{0pt}\setlength{\leftskip}{2em}}
  {}          % Indent amount
  {\bfseries} % Theorem head font
  {}          % Punctuation after theorem head
  { }         % Space after theorem head
  {\thmnumber{#2}}

\theoremstyle{customstyle}
\newtheorem{assumption}{}
\renewcommand{\theassumption}{(A\arabic{assumption})} 

\providecommand{\assumptionautorefname}{Assumption}
\newcommand{\lemmaautorefname}{Lemma}
\newcommand{\remarkautorefname}{Remark}
\newcommand{\corollaryautorefname}{Corollary}
\newcommand{\definitionautorefname}{Definition}
\newcommand{\propositionautorefname}{Proposition}
\renewcommand{\sectionautorefname}{Section}
\renewcommand{\subsectionautorefname}{Subsection}

\theoremstyle{definition}
\newtheorem{definition}{Definition}[section]

\theoremstyle{remark}
\newtheorem{remark}{Remark}[section]

\newcommand{\LM}[1]{\textcolor{red}{#1}}


\begin{document}

\maketitle

%\tableofcontents

\section{Linear Operators}

\begin{theorem}[Theorem 6.29 (Chapter III) of~\cite{kato2013perturbation}]
    If $T$ is closed in $\Space$ such that for some $\zeta$, $R(\zeta, T)$ exists and is compact, then the spectrum of $T$ consists of isolated eigenvalues (discrete spectrum) with finite multiplicities, and $R(\zeta)$ is compact for every $\zeta \in P(T)$.
    Here $P(T)$ indicates the resolvent set, where the resolvent operator is well-defined and bounded.  
\end{theorem}

\begin{theorem}[Theorem 1.16 (Chapter IV) of~\cite{kato2013perturbation}]
    Let $T \colon \D(T)\subset X \longrightarrow X$ be a closed linear operator whose inverse \(T^{-1}\in\B(X)\) is bounded on all of \(X\).
    Let $A \colon \D(A)\longrightarrow X$ be another linear operator with \(\D(A) \supset \D(T)\), and suppose there exist nonnegative constants \(a,b\) such that
    \[
    \|A\,x\| \le a\,\|x\| + b\,\|T\,x\|, \qquad \forall\,x\in\D(T),
    \]
    and moreover
    \[
    a\,\|T^{-1}\| + b <1.
    \]
    Define the perturbed operator $S := T + A$. 
    Then \(S\) is closed and bijective \(\D(S)\to X\), its inverse satisfies the norm-estimate
    \[
        \|S^{-1}\| \le \frac{\|T^{-1}\|}{1 - \bigl(a\,\|T^{-1}\|+b\bigr)}.
    \]
   If, in addition, \(T^{-1}\) is compact on \(X\), then so is \(S^{-1}\).
\end{theorem}


\begin{theorem}[Theorem 3.17 (Chapter IV) of~\cite{kato2013perturbation}]
    Let $T \colon \D(T)\subset X \longrightarrow X$ be a closed linear operator.
    Let $A \colon \D(A)\longrightarrow X$ be a $T$-bounded linear operator with \(\D(A) \supset \D(T)\).
    Let
    \[
        P(T) = \bigl\{\zeta\in\C : T - \zeta I \hbox{ is bijective and }(T-\zeta I)^{-1}\in\B(X)\bigr\}
    \]
    be the resolvent set of \(T\), and for each \(\zeta\in P(T)\) write $R(\zeta,T) :=(T-\zeta I)^{-1}$.
    Assume there exists some \(\zeta\in P(T)\) such that
    \[
        a\,\bigl\|R(\zeta,T)\bigr\| + b\,\bigl\|T\,R(\zeta,T)\bigr\| <1.
    \]
    Define the perturbed operator $S := T + A$. 
    Then \(\zeta\in\rho(S)\), \(S\) is closed and \((S-\zeta I)^{-1}\) exists and is bounded.
    The resolvent \(R(\zeta,S)=(S-\zeta I)^{-1}\) satisfies
    \[
        \bigl\|R(\zeta,S)\bigr\| \le \frac{\|R(\zeta,T)\|}{1 - \bigl(a\,\|R(\zeta,T)\| + b\,\|T\,R(\zeta,T)\|\bigr)}.
    \]
    If \(R(\zeta,T)\) is compact on \(X\), then so is \(R(\zeta,S)\). Hence \(S\) has compact resolvent.
\end{theorem}

\begin{theorem}[Theorem 6.16 (Chapter III) of~\cite{kato2013perturbation}]
    Let $T$ be a closed invertible operator in $\Space$.
    $\tilde{\Sigma}(T)$ and $\tilde{\Sigma}(T^{-1})$ are mapped onto each other by the mapping $\zeta \mapsto \zeta^{-1}$ of the extended complex plane. 
    This is a special case of the spectral mapping theorem, that states that the spectrum of $\phi(T)$ is the image under $\phi$ of the spectrum of $T$.
\end{theorem}

\begin{theorem}[Theorem XIII.64 of~\cite{reed1972methodsIV}]
    Let \( A \) be a self-adjoint operator that is bounded from below. Then the following are equivalent:
    \begin{enumerate}
        \item \( (A - \mu)^{-1} \) is compact for some \( \mu \in \rho(A) \).
        \item \( (A - \mu)^{-1} \) is compact for all \( \mu \in \rho(A) \).
        \item \( \{ \psi \in \mathcal{D}(A) : \|\psi\| \leq 1,\ \|A\psi\| \leq b \} \) is compact for all \( b \).
        \item \( \{ \psi \in \mathcal{Q}(A) : \|\psi\| \leq 1,\ \langle \psi, A\psi \rangle \leq b \} \) is compact for all \( b \).
        \item There exists a complete orthonormal basis \( \{ \varphi_n \}_{n=1}^\infty \subset \mathcal{D}(A) \) such that
        \[
        A \varphi_n = \mu_n \varphi_n,\quad \mu_1 \leq \mu_2 \leq \cdots,\quad \text{and } \mu_n \to \infty.
        \]
        \item $\mu_n(A) \to \infty$, where 
        \[
        \mu_n(A) \coloneqq \min_{S \subset D(A), |S|=n} \max_{0 \neq \psi \in S} \frac{\langle \psi, A\psi \rangle}{\|\psi\|^2}.
        \]
    \end{enumerate}
\end{theorem}


\section{Schrodinger Operator}

\begin{theorem}[Theorem XIII.67 of~\cite{reed1972methodsIV}]
    Let $V \in L_{\mathrm{loc}}^1$ be bounded from below and suppose that $V \to \infty$.
    Then $H = -\Delta + V$ defined as a aum of quadratic forms is an operator with compact resolvent.
    In particular, $H$ has a purely discrete spectrum and a complete set of eigenfunctions.
\end{theorem}


\bibliographystyle{plainnat}
\bibliography{biblio}

\end{document}